% 第7章 总结与展望
\chapter{总结与展望}
\label{chap:conclusion}

\section{本文工作总结}
\label{sec:summary}

本文围绕阿秒电子显微术的核心问题——如何在常规透射电子显微镜中实现亚周期时间分辨的电磁近场成像——开展了系统的研究。主要工作包括：

（1）建立了基于连续波激光驱动的阿秒电子脉冲产生方案。利用有质动力相互作用和时间Talbot效应，在电子波函数中刻印阿秒级时间结构，产生了$\sim$\SI{800}{as}的电子脉冲串。

（2）开发了阿秒场周期对比度成像方法。通过能量滤波器选择光子诱导能量增益的电子，实现了对纳米结构近场的空间分辨和时间分辨同时测量。

（3）在三类纳米光子系统中开展了亚周期动力学研究：针尖结构的手性表面波、介质共振器的多极模式干涉、多天线超原子的对称性破缺。

\section{主要创新点}
\label{sec:innovation}

本文的主要创新点如下：

\begin{enumerate}
    \item 提出了利用连续波激光在常规TEM中产生阿秒电子脉冲的实验方案，无需飞秒激光系统即可实现$\sim$\SI{800}{as}时间分辨率。
    \item 在亚周期时间尺度上观测到了手性表面波的方向性传播，揭示了自旋-轨道耦合在纳米光子激发中的矢量特性。
    \item 直接测量了介质纳米共振器中偶极和四极Mie模式的相对相位延迟，为Huygens超表面的时域设计提供了实验依据。
    \item 定量表征了$\sim$\SI{10}{nm}加工偏差对多天线系统电磁响应对称性的影响，为超表面的制造精度要求提供了参考。
\end{enumerate}

\section{展望}
\label{sec:outlook}

阿秒电子显微术作为一个新兴的研究方向，具有广阔的发展空间。未来的研究方向包括：

（1）\textbf{非重复动力学的单次成像}。目前的CW激光方案仅适用于周期性动力学过程。通过单周期束控技术\cite{Baum2007}或三波有质动力偏转方案\cite{Morimoto2020}，有望实现对非重复事件（如相变、随机开关）的阿秒时间分辨成像，代价是信噪比降低$\sim$100倍。

（2）\textbf{空间分辨率的提升}。当前$\sim$\SI{50}{nm}的空间分辨率受能量滤波器狭缝宽度限制。通过像差校正光学系统\cite{Morimoto2018}，有望将空间分辨率推进至$<\SI{10}{nm}$。

（3）\textbf{矢量场重构}。目前的技术测量的是电场沿电子轨迹的投影，而非完整的三维矢量场。通过倾转系列断层扫描方法，可以重构矢量场分布，代价是采集时间增加$\sim$10倍。

（4）\textbf{应用领域的拓展}。阿秒电子显微术有望应用于拓扑光子结构的近场表征、二维材料中的超快载流子动力学、以及量子材料中的相干声子激发等前沿问题。
